\documentclass[11pt]{article}
\usepackage{amsmath,amssymb,amsthm}
\title{Two Nested Certificates in the Antisymmetric Root Family}
\author{}\date{}
\begin{document}\maketitle
\paragraph{Status.}
These are complete exact computational certificates for two specified
butterfly targets, at sixteen and thirty-two wires. They prove feasibility
for every sufficiently small positive parameter. They are not an
all-target theorem, a density bound, a Lean theorem, or an asymptotic
head-complexity improvement. All original model definitions are unchanged.

\paragraph{The restricted family.}
Use the butterfly ordering, poles, and reflected queries of
\texttt{butterfly-adaptive-haar-biased.tex}. Let $H=2^d$, $q=H/2$,
\[
 p_y=\sum_{j=0}^{d-1}y_j\epsilon^{2^j},\quad
 z_y=p_y(1-p_y),\quad Q(t)=\prod_{y=0}^{H-1}(t-p_y).
\]
For the specified target permutation $\pi$, let
$u_y=1$ if $\pi^{-1}(y)<q$ and $u_y=-1$ otherwise. Restrict the residues to
\[
 c_y=\frac{(p_y-1/2)R(z_y)}{u_yQ'(p_y)},\qquad \deg R<q.       \tag{1}
\]
The one-node theorem in \texttt{butterfly-positive-root.tex} proves that
strict positivity of these residues automatically gives every root sign
at $1-p_s$, without a root bias. The remaining requirements are residue
positivity and the mixed-label comparisons in every nonroot block.
They are homogeneous linear inequalities in $q$ parameters.

For numerical conditioning only, parameterize $R$ by its values at
$z_0,\ldots,z_{q-1}$. Lagrange interpolation then makes (1) an explicit
$H$ by $q$ rational matrix $V(\epsilon)$. The certificate row list starts
with the $H$ positivity rows of $V$, followed by nonroot comparisons in
switch-block order, with positive-prefix index first and negative-prefix
index second. Each comparison is precisely formula (2) of the biased
Haar note. The JSON files specify $q$ rows whose equalities to one form a
square basis.

\paragraph{Exact finite and continuation checks.}
The files \texttt{butterfly-positive-root-family-h16.json} and
\texttt{butterfly-positive-root-family-h32.json} contain the targets,
exponent lists, and active row indices. The solver-free checker
\texttt{butterfly-positive-root-family-check.py} reconstructs all rows.
In finite mode it inverts the rational basis at $\epsilon=1/100$, checks
all inequalities with right side one, and independently checks every
root sign by the reciprocal sum. It does not rely on a solver's
feasibility report.

With \texttt{--laurent}, it uses exact rational series at tracked
precision $512$. It multiplies $V$ by the positive monomial
$\epsilon^E$, where
\[
 E=\sum_{j=0}^{d-1}2^{d-1-j}2^j=d2^{d-1}.
\]
This is the valuation of every $Q'(p_y)$, so all matrix entries become
analytic. It multiplies the basis right side by $\epsilon^{128}$ to
accommodate Laurent poles in its solution. Gaussian elimination pivots
by least valuation. Every division checks that its numerator precision
and valuation suffice; unknown zero multipliers are never skipped.
The checker substitutes into the original basis equations and checks
that every required score has a strictly positive leading term before
its tracked precision ends. Nonzero pivots give an invertible
rational-function basis. Consequently the computed leading terms are
those of actual rational functions, and positivity holds simultaneously
on some punctured real interval.

At sixteen wires there are $25$ inequalities, including $16$ residue
positivities. The eight pivot valuations are
\[
 0,0,0,0,0,2,4,6.
\]
The largest residue pole order is $6$; each residue leading coefficient
is $1$ or $2$. At thirty-two wires there are $73$ inequalities, including
$32$ residue positivities. The sixteen pivot valuations are eleven
zeros followed by
\[
 2,6,8,8,18.
\]
The largest residue pole order is $14$; all residue leading coefficients
belong to $\{1,4,8,9\}$. The JSON file records these leading terms as an
additional independently checked assertion. The smallest remaining
precision after removal of the right-side shift is $262$.
The original targets have respectively $32$ and $80$ query signs.
Their nonroot block biases are reconstructed using the strict gap
between the largest negative-label score and the smallest positive-label
score. Thus these finite-dimensional certificates establish all their
query signs, not only a selected subset.

\paragraph{The remaining uniform problem.}
The root polynomial has only half as many parameters as unrestricted
residues, yet both saved targets remain feasible and continue to
arbitrarily compact pole configurations. This is useful evidence about
one structured family. It does not prove that every target is feasible.
In a strict coalesced parent split, its leading positive residues form
a common point of the two child target cones. The one-node root theorem
does not supply that intersection. Pairwise compatibility of growing
child cones, or another way to avoid it, remains unproved.
\end{document}
